\section{Homotopy exact sequences}\label{XIII.4}\setcounter{subsection}{-1}
\subsection{}
\label{XIII.4.0}
Let $X$ and $S$ be two connected schemes, $f\colon X\to S$ a
morphism, $a$ a geometric point of~$X$, $\LL$ a set
of prime numbers. Let $K$ be the kernel of the canonical homomorphism
$\pi_1(X,a)\to\pi_1(S,a)$ and $N$ the smallest normal
pro-subgroup of~$K$ such that $K/N$ is a pro-$\LL$-group
$K^{\LL}$. Then $N$ is normal in $\pi_1(X,a)$ and we denote by
$$
\pi_1'(X,a)
$$
the
% original p. 415
quotient of~$\pi_1(X,a)$ by $N$. If $a$ is a geometric
point of a geometric fiber~$X_{\bar{s}}$, the
canonical morphisms
$$
\pi_1(X_{\bar{s}},a)\to \pi_1(X,a) \to
\pi_1(S,a)
$$
allow one to define canonical morphisms
$$
\pi_1^{\LL}(X_{\bar{s}},a)\lto{u} \pi_1'(X,a) \lto{v} \pi_1(S,a)
$$
We have $vu=0$.

\begin{proposition}
\label{XIII.4.1}
Let $S$ be a connected scheme, $f\colon X\to S$ a
locally $0$-acyclic morphism \textup{(SGA~4 XV~1.11)}; suppose moreover
that $f$ is $0$-acyclic (which, when $f$ is coherent, amounts to
saying that the geometric fibers of~$f$ are connected \textup{(SGA~4
XV~1.16)}). Let $\LL$ be a set of prime numbers. If $S'$ is a
scheme \'etale over~$S$, we denote by $X'$, $f'$ the inverse images
of~$X$, $f$ over~$S'$. Suppose that, for every \'etale covering
$S'$ of~$S$ and for every \'etale covering $E$ of
$X'$ that is a quotient of a Galois covering whose group is an
$\LL$-group, $(E,f')$ is cohomologically proper relative
to~$S'$ in dimension $\leqslant 0$ and $f'_* E$ is
constructible. Then, if $\bar{s}$ is a geometric point of
$S$ and $a$ a geometric point of the fiber $X_{\bar{s}}$, the
sequence of group homomorphisms
\begin{equation*}
\label{eq:XIII.4.1.1}
\tag{\thesubsection.1}
\pi_1^{\LL}(X_{\bar{s}},a)\lto{u} \pi_1'(X,a)\lto{v}\pi_1(S,a)\to1
\end{equation*}
is exact
\end{proposition}

This statement generalizes X.\Ref{X.1.4}, whose proof we shall
copy.

Let us first show that $v$ is surjective. It suffices to show that, for
every connected \'etale covering $S'$ of~$S$, $X'$ is also
connected (V~\Ref{V.6.9}). Let $C$ be a set having at least two
elements. It follows from the fact that $f$ is $0$-acyclic that
the canonical morphism
$$
\H^0(S',C_{S'})\to \H^0(X',C_{X'})
$$
is bijective, hence that $X'$ is connected, whence the surjectivity
of~$v$.

By
% original p. 416
definition of~$K^{\LL}$ \eqref{XIII.4.0}, we have the exact sequence
$$
1\to K^{\LL} \to \pi_1'(X,a)\to \pi_1(S,a)\to 1
$$
Let $\widetilde{S}$ be the universal covering of~$S$ and
$\widetilde{X}=\widetilde{S}\times_SX$; the group $K^{\LL}$ classifies the
Galois coverings $P$ whose group is an $\LL$-group, such that there
exist an \'etale covering $S'$ of~$S$ and a Galois covering
$Q$ of~$X'=X\times_SS'$ such that one has an isomorphism
$P\simeq Q\times_{X'}\widetilde{X}$. For the sequence
\eqref{eq:XIII.4.1.1} to be exact, it is necessary and sufficient that the
canonical morphism
$$
\pi_1^{\LL}(X_{\bar{s}},a)\to K^{\LL}
$$
be surjective. By the interpretation of~$K^{\LL}$ this
amounts to saying that, for every \'etale covering $S'$ of~$S$
and for every Galois covering $Q$ of~$X'$ whose group is an
$\LL$-group, such that $P=Q\times_{X'}\widetilde{X}$ is connected,
$Q|X_{\bar{s}}$ is then connected. Let us show that this last
condition is satisfied. Indeed, let $S'$ be an \'etale
covering of~$S$, $Q$ a Galois covering of~$X'$ whose group is an
$\LL$-group $F$, such that $Q|X_{\bar{s}}$ is disconnected, and let us show
that, after replacing $S'$ by an \'etale covering, $Q$
becomes disconnected. There exist a subgroup $G$ of~$F$ distinct from
$F$ and a torsor $R$ under $G|X_{\bar{s}}$ such that $Q|X_{\bar{s}}$
is obtained by extension of the structure group $G\to F$ from
$R$. The \'etale covering $E=Q/G$ of~$X'$ is such that
$E|X_{\bar{s}}$ has a section. By~\Ref{XIII.1.16} $f_*E$
is locally constant constructible and, after replacing $S'$
by an \'etale covering, one may even assume that $f_*E$
is constant. Since $(E,f')$ is cohomologically proper relative
to~$S'$ in dimension $\leqslant0$ and since $\H^0(X_{\bar{s}},
E|X_{\bar{s}})$ is nonempty, we see that $E$ has a section. But this
proves that $Q$ is disconnected, which completes the proof.

From~\Ref{XIII.4.1} we deduce the following lemma, which will be used
in~\Ref{XIII.4.6}.

\begin{lemma}
\label{XIII.4.2}
Let $S$ be a connected scheme, $f\colon X\to S$ a
locally $0$-acyclic and $0$-acyclic morphism, $\LL$ a set of
prime numbers. Suppose
% original p. 417
that, for every finite constant sheaf of~$\LL$-groups $F$ on~$X$,
$(F,f)$ is cohomologically proper relative to~$S$ in
dimension $\leqslant1$, and that, for every \'etale covering
$S'$ of~$S$ and for every \'etale covering $E$ of~$X'$
that is a quotient of a Galois covering whose group is an $\LL$-group,
$f'_*E$ is constructible. Then, if $\bar{s}$ is a
geometric point of~$S$ and $a$ a geometric point of the
fiber $X_{\bar{s}}$, the sequence of group homomorphisms
$$
\pi_1^{\LL}(X_{\bar{s}},a)\to\pi_1'(X,a)\to\pi_1(S,a)\to1
$$
is exact.
\end{lemma}

the hypotheses of~\Ref{XIII.4.1} are satisfied. It follows
indeed from~\Ref{XIII.1.13} 3) that, for every scheme $S'$
\'etale over~$S$ and for every \'etale covering $E$ of~$X'$
that is a quotient of a Galois covering whose group is an $\LL$-group,
$(E,f')$ is cohomologically proper relative to~$S'$ in
dimension $\leqslant 0$.
\begin{proposition}
\label{XIII.4.3}
Let $S$ be a connected scheme, $\LL$ a set of prime
numbers, $f\colon X\to S$ a $0$-acyclic morphism, locally
$1$-aspherical for $\LL$ \textup{(SGA~4 XV~1.11)}, $g\colon S\to X$ a
section of~$f$. Let $\bar{s}$ be a geometric point of~$S$,
$a$ a geometric point of the fiber $X_{\bar{s}}$. Suppose
that, for every constant sheaf of
$\LL$-groups~$F$, $(F,f)$ is cohomologically proper in dimension
$\leqslant1$, that the direct image under~$f$ of the stack of torsors under
$F$ is a $1$-constructible stack \eqref{XIII.0} and that, for every
\'etale covering $E$ of~$X'$ that is a quotient of a Galois
covering whose group is an $\LL$-group, $f'_*E$ is constructible. Then
the sequence of group homomorphisms
\begin{equation*}
\label{eq:XIII.4.3.1}
\tag{\thesubsection.1}
1\to\pi_1^{\LL}(X_{\bar{s}},a)\lto{u} \pi_1'(X,a)\lto{v}\pi_1(S,a)\to1
\end{equation*}
is exact.
\end{proposition}

Taking~\Ref{XIII.4.2} into account, it suffices to show the injectivity
of the morphism $u$, \ie to prove that, for every principal
covering $\overline{Z}$ of~$X_{\bar{s}}$
% original p. 418
whose group is an $\LL$-group~$C$, there exist an \'etale covering
$Z$ of~$X$ and a morphism from a connected component of~$Z|X_{\bar{s}}$
to~$\overline{Z}$ (V~\Ref{V.6.8}). So let $\overline{Z}$ be a
principal covering of~$X_{\bar{s}}$ whose group is an $\LL$-group~$C$, and~$\bar{z}$ its class in
$\H^1(X_{\bar{s}},C_{X_{\bar{s}}})$. By \Ref{XIII.1.5}~d), we
have a canonical isomorphism
$$
\left(\R^1f_*C_X\right)_{\bar{s}}\isomto\H^1(X_{\bar{s}},C_{X_{\bar{s}}}),
$$
and by~\Ref{XIII.1.16}, $\R^1f_*C_X$ is a locally constant
constructible sheaf. One can therefore find an
\'etale covering $S'$ of~$S$ such that $\R^1f_*C_X|S'$ is
constant. If $\bar{s}\to S'$ is a geometric point
above the geometric point $\bar{s}\to S$, there exists an
element~$z$ of~$\H^0(S', \R^1f_*C_X)$ whose image in
$\H^1(X_{\bar{s}},C_{X_{\bar{s}}})$ is $\bar{z}$. By
Lemma~\Ref{XIII.4.3.1} below, one can find an \'etale
covering $S'_1$ of~$S'$ and a torsor $P$ on~$X'_1$ with group~$C$
whose image in $\H^0(S'_1,\R^1f_*C)$ is equal to the
restriction of~$z$. The torsor~$P$ is representable by an
\'etale covering $Z$ of~$X'_1=X\times_SS'_1$ such that
$Z\times_{X'_1}X_{\bar{s}}$ is isomorphic to $\overline{Z}$. If
one regards $Z$ as an \'etale covering of~$X$, one
then has a morphism from~$Z\times_XX_{\bar{s}}$ to $\overline{Z}$, which
completes the proof.

\begin{sublemma}
\label{XIII.4.3.1}
Let $f\colon X\to S$ be a $0$-acyclic and locally
$0$-acyclic morphism, $g$ a section of~$f$. Let $C$ be a finite
constant group such that $(C_X,f)$ is cohomologically proper in dimension
$\leqslant0$ and such that the direct image under $f$ of the stack of torsors under
$C_X$ is constructible. Then, for every section $z$ of
$\H^0(S,\R^1f_*C_X)$, one can find an \'etale covering
$S_1$ of~$S$ and, if $X_1=X\times_SS_1$, an element of
$\H^1(X_1,C_{X_1})$ whose image under the canonical morphism
$$
\H^1(X_1,C_{X_1})\to\H^0(S_1,\R^1f_*C_{X_1})
$$
is equal to the restriction of~$z$ to
$\H^0(S_1,\R^1f_*C_{X_1})$.
\end{sublemma}

For every scheme $S'$ \'etale over~$S$, we set $X'=X\times_SS'$,
and we denote by $g'$ (\resp $F'$, etc.) the inverse image of~$g$ (\resp $F$,
etc.) under the morphism
% original p. 419
$S'\to S$. The presheaf $G$ on~$S$ defined by
\newlength{\tmpRacinet} \setlength{\tmpRacinet}{\textwidth}
\addtolength{\tmpRacinet}{-7em}
$$
G(S') = \left\{\text{\parbox{\tmpRacinet}{classes mod.\ isomorphism of
		torsors $P$ on~$X'$ with group $C_{X'}$, endowed with an
		isomorphism $g'^*P\isomto C_{S'}$ }} \right\}
$$
is then a sheaf. Indeed, this follows by descent from the fact
that an isomorphism of a torsor $P$ on~$X'$ is completely
determined by its restriction to $g'(S')$. Moreover, we have a
surjective morphism
$$
G\to \R^1f_*F.
$$
Let $z$ be an element of~$\H^0(S,\R^1f_*C_X)$ and $H$ the
subsheaf of~$G$ that is the inverse image of~$z$. It suffices to show
that $H$ is a locally constant constructible sheaf. Now, this
property being local on~$S$, we may assume that $z$
comes from an element of~$\H^1(X,C_X)$ represented by
a torsor $P$ such that $g^*P$ is isomorphic to $C_S$. Giving an
isomorphism $i\colon g^*P\isomto C_S$ amounts to giving a
global section of~$\SheafAut_{C_S}(g^*P)$, and two isomorphisms $i$
and $i'$ define the same element of~$G(X)$ if and
only if $ii'^{-1}$ is the image of an element of
$\Aut_{C_X}(P)$. If one considers the canonical injection
$$
f_*\SheafAut_{C_X}(P)\to\SheafAut_{C_S}(g^*P)\simeq C,
$$
$H$ is therefore identified with the quotient of~$\SheafAut_{C_S}(g^*P)$ by
$f_*\SheafAut_{C_X}(P)$. By~\Ref{XIII.1.16}
$f_*\SheafAut_{C_X}(P)$ is locally constant; hence the same
holds for~$H$, which completes the proof.

\begin{examples}
\label{XIII.4.4}
Let us note that, if $S$ is a connected scheme, the hypotheses
of~\Ref{XIII.4.1} are satisfied when the morphism $f$ is proper,
flat, of finite presentation, with separable connected geometric
fibers, $\LL$ being arbitrary
(\cf X~\Ref{X.1.3}). The hypotheses of~\Ref{XIII.4.3} are
satisfied if moreover $f$ is smooth and has a section, if one
denotes by $\LL$ the set of prime numbers distinct from the
residue characteristics of~$S$ (SGA~4 XV~2.1 and XVI~5.2).

The hypotheses of~\Ref{XIII.4.1} are also satisfied if $S$ is connected,
if one has a scheme~$Z$ proper of finite presentation, flat over~$S$, with
% original p. 420
separable connected geometric fibers, such that $X$ is
the complement in $Z$ of a divisor with normal crossings
relative to~$S$, $\LL$ being the set of prime
numbers distinct from the residue characteristics of~$S$
\eqref{XIII.2.9}. The hypotheses of~\Ref{XIII.4.3} are
satisfied if moreover $f$ is smooth and has a section.
\end{examples}

\subsection{}
\label{XIII.4.5}
Let us take up again the notation and hypotheses of~\Ref{XIII.4.3}. If
$\bar{s}$ is a geometric point of~$S$ and $a=g(\bar{s})$,
the section $g$ allows one to define a morphism
$$
w\colon \pi_1(S,a)\to \pi_1'(X,a),
$$
so that $\pi_1'(X,a)$ is identified with the semi-direct product of
$\pi_1(S,a)$ by $\pi_1(X_{\bar{s}},a)$. The profinite group
$\pi_1(S,a)$ therefore acts on~$\pi_1(X_{\bar{s}}, a)$. Since
$\pi_1^\LL(X_{\bar{s}},a)$ is a strict projective limit of groups
invariant under the action of~$\pi_1(S,a)$, the datum of
$\pi_1^\LL(X_{\bar{s}},a)$ endowed with this action is equivalent
to the datum of a strict projective system of finite \'etale
group schemes over~$S$, which we denote by
$$
\pi_1^\LL(X/S,g,\bar{s})\quad\text{or
simply}\quad\pi_1^\LL(X/S,g).
$$

We then have the following properties:
\setcounter{subsubsection}{0}

\subsubsection{}
\label{XIII.4.5.1}
For every finite \'etale group scheme $G$ over~$S$ whose
fibers are $\LL$-groups, the set $E$ of classes of torsors
$P$ under the inverse image $G_X$ of~$G$ on~$X$, endowed with an isomorphism
$g^*P\isomto G$, is canonically isomorphic to the set
$$
\Hom_S(\pi_1^\LL(X/S,g,\bar{s}),G) \quad\text{mod.\ inner
automorphisms of~$G$.}
$$

\subsubsection{}
\label{XIII.4.5.2}
For every finite \'etale group scheme $G$ over~$S$ whose
fibers are $\LL$-groups, the sheaf $\R^1f_*G_X$ is
canonically isomorphic to the sheaf associated with the presheaf
$$
S'\mto \Hom_{S'}(\pi_1^\LL(X/S,g,\bar{s}),G)\
\text{mod.\ inner automorphisms of~$G$.}
$$
($S'$ denotes a scheme \'etale over~$S$).

\subsubsection{}
\label{XIII.4.5.3}
Let
% original p. 421
$S'$ be a connected
$S$-scheme, $\bar{s}$ a geometric point of~$S'$, $X'$, $g'$
the respective inverse images of~$X,g$ over~$S'$. Then
$\pi_1^\LL(X'/S',g',\bar{s})$ is canonically isomorphic to
the inverse image of~$\pi_1^\LL(X/S,g,\bar{s})$ over~$S'$. For every
geometric point $\xi$ of~$S$, the fiber
$\pi_1^\LL(X/S,g,\bar{s})_\xi$ is isomorphic to $\pi_1^\LL(X_\xi)$.

Indeed, the datum of~$G$ is equivalent to the datum
of an abstract $\LL$-group $\mathbf{G}$
on which $\pi_1(S,a)$ acts, whence an action of
$\pi_1'(X,a)$ on~$\mathbf{G}$. The isomorphism defined
in~\Ref{XIII.4.5.1} is then obtained by restriction to the subset
$E$ from the canonical morphism
\[
\H^1(\pi_1'(X,a),\mathbf{G})\to
\H^1(\pi_1^\LL(X_{\bar{s}},a),\mathbf{G}) =
\Hom(\pi_1^\LL(X_{\bar{s}},a),\mathbf{G})/\text{inn.\ aut.\ }G,
\]
the set $E$ being mapped bijectively onto the subset of
morphisms from $\pi_1^\LL(X_{\bar{s}},a)$ to $\mathbf{G}$ that are
compatible with the action of
$\pi_1(S,a)$. Assertion~\Ref{XIII.4.5.3} then follows from the
definition of~$\pi_1^\LL(X/S,g,\bar{s})$ taking into account the homotopy
exact sequence \eqref{eq:XIII.4.3.1}, and~\Ref{XIII.4.5.2} is
deduced from \Ref{XIII.4.5.1} and~\Ref{XIII.4.5.3}.

\begin{proposition}[K\"unneth formula]
\label{XIII.4.6}
Let $k$ be a separably closed field of characteristic
$p\geq 0$, $X$ and $Y$ two connected $k$-schemes, $a$ a
geometric point of~$X$, $b$ a geometric point of~$Y$,
$c$ a geometric point of~$X\times_k Y$ above~$a$ and
$b$. Suppose that one of the following two conditions is satisfied:
\begin{enumerate}
\item[a)] $X$ is of finite type over~$k$ and the schemes of finite type
over an algebraic closure $\bar{k}$ of~$k$, of dimension
$\leq\dim X$, are strongly desingularizable. \textup{(SGA~5 I~3.1.5)}.
\item[b)] $X$ is quasi-compact and quasi-separated and every
scheme of finite type over~$\bar{k}$ is strongly
desingularizable.
\end{enumerate}

Then, if $p'$ is the set of prime numbers distinct from~$p$,
the morphism
\begin{equation*}
\label{eq:XIII.4.6.0} \tag{\thesubsection.0} {}
\pi_1^{p'}(X\times_k Y,c)\to\pi_1^{p'}(X,a)\times\pi_1^{p'}(Y,b),
\end{equation*}
deduced
% original p. 422
from the homomorphisms on the fundamental groups associated with the
projections
$$
X\times_k Y\to X \quad \text{and} \quad X\times_k Y\to Y,
$$
is an isomorphism.
\end{proposition}

We may assume $k$ algebraically closed and $X$ reduced
(SGA~4 VIII~1.1). Let $Z=X\times_k Y$, and let $g\colon X\to k$ and $f\colon
Z\to Y$ be the canonical morphisms. The morphism $g$, hence also $f$,
is universally locally $1$-aspherical for $p'$
by~\Ref{XIII.3.5}. Since $X$ is connected, $f$ is
$0$-acyclic (SGA~4 XV~1.16). On the other hand it follows
from~\Ref{XIII.3.2} that, for every finite $p'$-group $C$, $(C_X,g)$ is
cohomologically proper relative to $k$ in dimension $\leq
1$. It follows that $(C_Z,f)$ is cohomologically proper
relative to $Y$ in dimension $\leq 1$ (\Ref{XIII.1.5}~c)) and
that $f_*C_Z$ and $\R^1f_*C_Z$ are constant sheaves. Consequently
$g$ satisfies all the hypotheses of~\Ref{XIII.4.2}. We
therefore have the exact sequence
$$
\pi_1^{p'}(X_b,c)\to\pi_1^{p'}(Z,c)\to\pi_1^{p'}(Y,b)\to 1.
$$
Moreover the composite morphism
$$
\pi_1^{p'}(X_b,c)\to\pi_1^{p'}(Z,c)\to\pi_1^{p'}(Y,b)
$$
is an isomorphism and we therefore have the exact sequence
$$
1\to\pi_1^{p'}(X,a)\to\pi_1^{p'}(Z,c)\to\pi_1^{p'}(Y,b)\to 1.
$$
On the other hand the morphism \eqref{eq:XIII.4.6.0} allows one to define
a morphism from this exact sequence to the exact sequence
$$
1\to\pi_1^{p'}(X,a)\to
\pi_1^{p'}(X,a)\times\pi_1^{p'}(Y,b)\times\pi_1^{p'}(Y,b)\to 1,
$$
and it follows that the morphism \eqref{eq:XIII.4.6.0} is an
isomorphism.

\subsection{}
\label{XIII.4.7} Let
$$
\xymatrix@C=1.2cm{
X \ar[d]_-f&\ar[l] X_U \ar[d]_-{f_U}& \ar[l]X_{\bar{s}}\ar[d] \\
S&\ar[l] U &\ar[l]\bar{s}
}
$$
be a
% original p. 423
diagram whose squares are cartesian, where $S$ is an
arcwise connected scheme (SGA~4 IX~2.12), $U$ a connected open subset of
$S$, $\bar{s}$ a geometric point of~$U$. Let $a$ be a
geometric point of~$X_{\bar{s}}$, $\LL$ a set of
prime numbers. Let $g$ be a section of~$f$ and suppose that the
following conditions are satisfied:
\begin{enumerate}
\item[a)] The morphism $f$ is $0$-acyclic, locally $0$-acyclic,
and, for every \'etale covering $S'$ of~$S$ and every
\'etale covering $E$ of~$X\times_S S'$ that is a quotient of a
Galois covering whose group is an $\LL$-group, $(E,f_{(S')})$ is
cohomologically proper relative to~$S'$ in dimension $\leq
0$.
\item[b)] The morphism $f_U$ is locally $1$-aspherical for
$\LL$, and, for every finite constant sheaf of~$\LL$-groups $F$ on~$X_U$, $(F,f_U)$ is cohomologically proper in dimension $\leq 1$ and
the fibers of~$\R^1 f_{U*}F$ are finite.
\end{enumerate}

From~\Ref{XIII.4.1} and~\Ref{XIII.4.3} one then deduces the following
commutative diagram, whose rows are exact:
\begin{equation*}
\label{eq:XIII.4.7.0} \tag{\thesubsection.0} {}
\begin{array}{c}
\xymatrix{
1 \ar[r]& \pi_1^{\LL}(X_{\bar{s}},a) \ar[r]\ar@{=}[d]& \pi'_1(X_U,a) \ar[r]\ar[d]& \pi_1(U,a)\ar[r]\ar[d]& 1 \\
&\pi_1^{\LL}(X_{\bar{s}},a) \ar[r]&\pi'_1(X,a) \ar[r]& \pi_1(S,a) \ar[r]& 1
}
\end{array}
\end{equation*}

Thanks to the section $g$, we have morphisms
$$\pi_1(U,a)\to\pi'_1(X_U,a),\quad \pi_1(S,a)\to\pi'_1(X,a)\,\;$$
from which one
deduces a morphism from the amalgamated sum of~$\pi_1(S,a)$ and
$\pi'_1(X_U,a)$ over~$\pi_1(U,a)$ to $\pi'_1(X,a)$:
\begin{equation*}
\label{eq:XIII.4.7.1} \tag{\thesubsection.1} {}
\varphi\colon\pi=\pi_1(S,a)\coprod_{\pi_1(U,a)}\pi'_1(X_U,a)\to\pi'_1(X,a).
\end{equation*}
Suppose that the following condition is satisfied:
\begin{enumerate}
\item[c)] If $T=S-U$, we have $\prof \et_T(S)\geq 2$ (SGA~2 XIV~1.1).
\end{enumerate}
Then the functor which associates with an \'etale covering of~$S$
its restriction to $U$ is fully faithful
(SGA~2 XVI~1.4). It follows that the morphism
$\pi_1(U,a)\to\pi_1(S,a)$ is surjective (V~\Ref{V.6.9}) and one
deduces from the
% original p. 424
diagram \eqref{eq:XIII.4.7.0} that the same holds for the morphism
$\pi'_1(X_U,a)\to\pi'_1(X,a)$; a fortiori $\varphi$ is an
epimorphism. Let
\begin{equation*}
\label{eq:XIII.4.7.2} \tag{\thesubsection.2} {}
K=\Ker(\pi_1(U,a)\to\pi_1(S,a)).
\end{equation*}
The group $\pi$ of \eqref{eq:XIII.4.7.1} is identified with the quotient of
$\pi'_1(X_U,a)$ by the closed invariant subgroup generated
by the image $L$ of~$K$ in $\pi'_1(X_U,a)$. Let us regard
$\pi'_1(X_U,a)$ as the semi-direct product of~$\pi_1(U,a)$ by
$\pi_1^\LL (X_{\bar{s}},a)$. The group $K$ then acts by
inner automorphisms on~$\pi_1^\LL (X_{\bar{s}},a)$, and the
quotient $\pi=\pi'_1(X_U,a)/L$ is identified with the semi-direct product of
$\pi_1(U,a)/K=\pi_1(S,a)$ by the group $\pi_1^\LL(X_{\bar{s}},a)_K$
of coinvariants of~$\pi_1^{\LL}(X_{\bar{s}},a)$ under~$K$. We
finally have an epimorphism
\begin{equation*}
\label{eq:XIII.4.7.3} \tag{\thesubsection.3} {}
\varphi\colon\pi=\pi_1^\LL(X_{\bar{s}},a)_K\cdot\pi_1(S,a)\to\pi'_1(X,a).
\end{equation*}

The proposition which follows gives conditions under which the
morphism $\varphi$ is an isomorphism.

\setcounter{subsubsection}{3}

\begin{subproposition}
\label{XIII.4.7.4} The notation is that of~\Ref{XIII.4.7}. We
suppose that, in addition to conditions \textup{a), b), c)}, the following
conditions are satisfied:
\begin{enumerate}
\item[d)] For every point $t$ of~$T=S-U$, the morphism $f$ is
locally $1$-aspherical for $\LL$ at $g(t)$.
\item[e)] For every point $t\in T$, every irreducible component
of the fiber $X_t$ contains $g(t)$ and, for every point $x$ of
$X_t-\{g(t)\}$ which is not maximal, we have
$$
\prof\hop_x(X)\geq 3 \quad\textup{(SGA~2 XIV~1.2)}
$$
and the ring $\cal{O}_{X,x}$ is noetherian.
\end{enumerate}

Then the morphism \eqref{eq:XIII.4.7.3} is an isomorphism.
\end{subproposition}

As was said previously, the group $\pi$ is identified with the
quotient of~$\pi'_1(X_U,a)$ by the closed invariant subgroup $L$
generated by the image of~$K$ \eqref{eq:XIII.4.7.2} in
$\pi'_1(X_U,a)$. This amounts to saying that $\pi$ classifies the
% original p. 425
principal coverings $Z$ of~$X_U$ such that $g^{-1}_U(Z)$
extends to an \'etale covering of~$S$, and which induce on~$X_{\bar{s}}$ a covering obtained by extension of the structure
group from a principal covering whose group is an
$\LL$-group. To prove that $\varphi$ is an isomorphism, it
suffices to show that such a covering $Z$ extends to all
of~$X$.

Let us first show that $Z$ extends to an \'etale covering
of an open set containing $X_U$ and $g(S)$. Let $W$ be a scheme
\'etale over~$X$ whose image contains $X_U$ and $g(S)$, and let us set
$W_U=W\times_S U$. From the fact that the morphism $W\to S$ is $0$-acyclic
and from the relation $\mathrm{prof}.\,\et_T S\geq 2$, it follows that we have
$$
\mathrm{prof}.\,\et_{(W-W_U)} W\geq 2
$$
(SGA~2 XIV~1.13); consequently, if $Z|W_U$ extends to an
\'etale covering of~$W$, this extension is unique up to
unique isomorphism. It follows that the problem of
extending $Z$ to a neighborhood of~$g(S)\cup X_U$ is local for the
\'etale topology in the neighborhood of the points of~$g(T)$. If $t$ is a
point of~$T$, we set $x=g(t)$, and we denote by $\bar{X}$ (\resp $\bar{S}$)
the strict localization of~$X$ at $\bar{x}$ (\resp of~$S$ at
$\bar{t}$), $\bar{U}=U\times_S \bar{S}$, $\bar{X}_U=\bar{X}\times_X
X_U$, $\bar{g}\colon\bar{S}\to\bar{X}$ the morphism deduced from
$g$. It suffices to show that, for every point $t$ of~$T$, the inverse
image $\bar{Z}$ of~$Z$ on~$\bar{X}_U$ extends to $\bar{X}$
or, what amounts to the same, is trivial. Now, by definition
of~$Z$, the inverse image of~$\bar{Z}$ on~$\bar{U}$ is
trivial. To prove that $\bar{Z}$ is trivial, it suffices to show
that it is of the form $\bar{f}^*_U E$, where $E$ is a
principal covering of~$\bar{U}$; indeed we shall then have
$E\simeq\bar{g}^*_U\bar{f}^*_U E \simeq\bar{g}^*_U\bar{Z}$, whence
the result since $\bar{g}^*_U\bar{Z}$ is trivial. Now, the
morphism $\bar{f}_U$ being $0$-acyclic and locally $0$-acyclic,
it suffices, in order to prove that $\bar{Z}$ comes from~$\bar{U}$, to
show that, for every geometric point algebraic over a
point of~$\bar{U}$, which one may assume to be the point
$\bar{s}$, $\bar{Z}|X_{\bar{s}}$ is trivial (SGA~4 XV~1.15). But since
$\bar{Z}|X_{\bar{s}}$ is obtained by extension of the structure
group from a principal covering whose group is an
$\LL$-group, this follows from the fact that the morphism $\bar{f}$ is
$1$-aspherical for $\LL$.

We
% original p. 426
have therefore shown that there exists an open neighborhood $V$ of
$g(S)\cup X_U$ such that $Z$ extends to an \'etale covering
$Z_V$ of~$V$. Let us show that $Z_V$ extends to all of~$X$.
It suffices to see that, for every point $x$ of~$X-V$, we have
$$
\prof\hop_x X\geq 3.
$$
Now this follows from hypothesis e) and from the fact that a point $x$
of~$X-V$ cannot be maximal in its fiber $X_t$ since, every
irreducible component of~$X_t$ containing $g(t)$, every maximal
point of~$X_t$ belongs to $V$.

\begin{corollary} \label{XIII.4.8}
The hypotheses are those of~\Ref{XIII.4.7.4}, but we assume moreover
that $\pi_1(S,a)=1$. Then we have an isomorphism
\begin{equation*}
\label{eq:XIII.4.8.*} \tag{$*$}
{}\pi_1^\LL(X,a)\isomto\pi_1^\LL(X_{\bar{s}},a)_K.
\end{equation*}
In particular, if $\pi_1^\LL(X_{\bar{s}},a)$ is topologically of
finite presentation and if $K$ acts on~$\pi_1^\LL(X_{\bar{s}},a)$ through a group of finite
type, then $\pi_1^\LL(X,a)$ is topologically of finite
presentation.
\end{corollary}

The isomorphism \eqref{eq:XIII.4.8.*} was proved
in~\Ref{XIII.4.7.4}. Suppose that $\pi_1^\LL(X_{\bar{s}},a)$ is the quotient
of the free pro-$\LL$-group on $n$ generators
$L(x_1,\dots,x_n)$ by the closed invariant subgroup generated
by the elements $y_1,\dots,y_p$ of~$L(x_1,\dots,x_n)$, and
suppose that $K$ acts through a group
generated by elements $k_1,\dots,k_q$. If for every
$i\in[1,n]$, $j\in[1,q]$, we denote by $z_{ij}$ an element of
$L(x_1,\dots,x_n)$ lifting the element $(k_j\cdot
x_i)x_i^{-1}$, then $\pi_1^\LL(X_{\bar{s}},a)_K$ is the quotient of
$L(x_1,\dots,x_n)$ by the closed invariant subgroup generated
by the elements $(y_i)_{i\in[1,p]}$,
$(z_{ij})_{i\in[1,n],j\in[1,q]}$.

\begin{remarks}
\label{rem:XIII.4.6}
a) Conditions a) to e) of~\Ref{XIII.4.7} are
satisfied when $S$ is a connected normal scheme, $U$ a
dense retrocompact open subset of~$S$, $f$ a proper morphism of
finite presentation, with fibers that are geometrically connected
and irreducible at every point $t$ of~$T$, $f$ being moreover
separable,
% original p. 427
smooth at the points of~$X_U\cup g(T)$, $\LL$ being the set of
prime numbers distinct from the residue characteristics of
$S$, and $X$ being regular at every point of~$X_t$. Condition
a) follows indeed from SGA~4 XV~4.1 and 1.4; conditions
b) and d) follow from~\Ref{XIII.1.4} and SGA~4 XV~2.1 and
XVI~5.2. Finally e) follows from SGA~XIV~1.11.
\par\smallskip
b) Corollary~\Ref{XIII.4.5} applies to compute the
fundamental group $\pi_1^{p'}(X)$ of a proper and smooth surface $X$
over a separably closed field $k$ of characteristic $p$
($p'$~denoting the set of prime numbers distinct from
$p$). The method was communicated to us by
J.P\ptbl \textsc{Murre}; it consists in reducing, by blowing up
$X$, to the case where one has a fibration $X\to\PP^1_k$ and an open set $U$
of~$\PP^1_k$ satisfying the hypotheses of~\Ref{XIII.4.7} (see
SGA~7 for more details). The same method can be
used more generally (\loccit) to prove that, if
$X$ is a connected $k$-scheme of finite type, and if the schemes
of finite type of dimension $\leq\dim X$ over an algebraic
closure of~$k$ are strongly desingularizable
(SGA~5 I~3.1.5), then $\pi_1^{p'}(X)$ is topologically of
finite presentation.
\end{remarks}
